\section{Glosario y herramientas}

\subsection{Consulta rápida de términos clave}

\begin{table}[H]
\centering
\begin{tabular}{lp{10cm}}
\toprule
Término & Definición \\
\midrule
TD error & $r + \gamma \max_{a'} Q(s', a') - Q(s, a)$, sirve para medir la magnitud de la actualización de la estimación actual de Q \\
Prioritized Experience Replay & Muestrea según el TD error e introduce la IS correction \\
Target Network & Red que se sincroniza periódicamente y estabiliza el TD target \\
State Action Coverage & Distribución de los distintos pares state-action en el replay buffer \\
\bottomrule
\end{tabular}
\caption{Términos principales de esta clase}
\end{table}

\begin{knowledgebox}{Inventario de la cadena de herramientas}
Entre las herramientas de uso común están: OpenAI Baselines (plantilla básica de DQN), RLlib (admite múltiples agents), MLFlow + TensorBoard (registro de métricas), WeTest o Supersonic (validación de la cuantización en el edge).
\end{knowledgebox}

\subsection{Resumen de la sección}

Dominar a la vez el glosario y la cadena de herramientas ayuda a pasar de la teoría a la práctica de ingeniería; en particular, evita repetir la exploración por ensayo y error al reproducir un benchmark y al poner un modelo en producción.

 \subsection{Resumen de la sección}

Dominar a la vez el glosario y la cadena de herramientas ayuda a pasar de la teoría a la práctica de ingeniería; en particular, evita repetir la exploración por ensayo y error al reproducir un benchmark y al poner un modelo en producción.

